\section*{Chapter \ref{ch:g12:monetary-policy} --- Monetary Policy}
% figdata: figdata/grade-12/10-monetary-policy/out/ (fedfunds, taylor, policy_key, brenland); Brenland invented.

\begin{solution}{exo:g12:monetary-policy:1}
In mid-1981 (June), at about 19\,\%.
\end{solution}

\begin{solution}{exo:g12:monetary-policy:2}
From the end of 2008 to the end of 2015, and from 2020 to early 2022.
\end{solution}

\begin{solution}{exo:g12:monetary-policy:3}
$2 + 3 + 0.5 \times 1 + 0.5 \times 1 = 6\,\%$.
\end{solution}

\begin{solution}{exo:g12:monetary-policy:4}
$2 + 1 + 0.5 \times (-1) + 0.5 \times (-2) = 1.5\,\%$.
\end{solution}

\begin{solution}{exo:g12:monetary-policy:5}
$5.25 - 3 = 2.25\,\%$.
\end{solution}

\begin{solution}{exo:g12:monetary-policy:6}
4, 7 and 10\,\%; real rates 2, 3 and 4\,\%: the real rate rises with inflation.
\end{solution}

\begin{solution}{exo:g12:monetary-policy:7}
$4 - 2 = 2$, $4 - 1 = 3$ and 4\,\%.
\end{solution}

\begin{solution}{exo:g12:monetary-policy:8}
Lower \omterm{def:g11:money-and-banking:policy-rate}{policy rate}; lower market rates for loans and \omterm{def:g12:present-value:bond}{bonds} and higher prices of shares and houses; more borrowing, investment and consumption, and a weaker currency
with higher \omterm{def:g10:national-accounts:net-exports}{net exports}; output rises towards or above potential; with less slack, wages and prices rise faster: inflation rises.
\end{solution}

\begin{solution}{exo:g12:monetary-policy:9}
The rate rises in June; inflation starts to fall because of it in June of the following year, about seventeen months after it began to rise. A bank acting on forecasts
would raise its rate when it expects inflation to rise, before the data show it.
\end{solution}

\begin{solution}{exo:g12:monetary-policy:10}
Rule 13.1\,\%, actual 1.68\,\%: a difference of about 11.4 points.
\end{solution}

\begin{solution}{exo:g12:monetary-policy:11}
The rate cannot go far below zero, since people can hold cash at zero. It can buy long-term \omterm{def:g12:present-value:bond}{bonds} (\omterm{def:g12:monetary-policy:qe}{quantitative easing}) and promise to keep its rate low for longer
(\omterm{def:g12:monetary-policy:qe}{forward guidance}).
\end{solution}

\begin{solution}{exo:g12:monetary-policy:12}
Neither, by itself: the rule describes how policy was set in some years, not what is best. A bank might depart from it because it expects the rise in inflation to be
temporary (supply shocks fading), because its estimate of the gap or of the neutral rate differs from the rule's inputs, or because of the lags: it sets policy for
the inflation it forecasts, not the inflation it sees.
\end{solution}

\begin{solution}{exo:g12:monetary-policy:13}
If wage setters expect 0\,\%, the bank gets 1 by letting 5\,\% happen and 0 by keeping 0\,\%: it lets 5\,\% happen, so expecting 0\,\% would be wrong. If they expect 5\,\%, the
bank gets $-1$ by letting 5\,\% happen and $-3$ by keeping 0\,\%: it lets 5\,\% happen, and the expectation is right, so this is an equilibrium. A bank that could commit
to 0\,\% would get 0, better than $-1$.
\end{solution}

\begin{solution}{exo:g12:monetary-policy:14}
Across sixteen rich countries, central banks under political control had higher and more variable inflation, with little difference in growth, unemployment or real
rates. It is a correlation: countries whose voters dislike inflation may both choose independent banks and keep inflation low by other means, so independence may not
be the cause.
\end{solution}

\begin{solution}{exo:g12:monetary-policy:15}
A: $2 + 4 + 1 + 1 = 8\,\%$. B: $2 + 0 - 1 - 1.5 = -0.5\,\%$. Union: inflation 2\,\%, gap $-0.5\,\%$: $2 + 2 + 0 - 0.25 = 3.75\,\%$, too low for A and too high for B.
\end{solution}

\begin{solution}{pb:g12:monetary-policy:1}
\textbf{1.} 2022: inflation 10.2\,\%, \omterm{def:g11:money-and-banking:policy-rate}{policy rate} 5.42\,\%.

\textbf{2.} 2019: $2.72 - 1.9 = 0.82\,\%$. 2022: $5.42 - 10.2 = -4.78\,\%$.

\textbf{3.} 2020, at $-4.3\,\%$: a sharp fall in demand and in production at once, such as a pandemic or a financial crisis.

\textbf{4.} Above: gaps of 0.9 to 2.2\,\%.

\textbf{5.} Yes: from 3.11\,\% to 5.42\,\%, by 2.31 points.

\textbf{6.} $2 + 0.9 + 0.5 \times (-1.1) + 0.5 \times 0.4 = 2.55\,\%$.

\textbf{7.} $2 + 2.2 + 0.1 - 2.15 = 2.15\,\%$.

\textbf{8.} $2 + 10.2 + 4.1 - 0.05 = 16.25\,\%$.

\textbf{9.} $16.25 - 5.42 = 10.83$ points.

\textbf{10.} $2 + 5.6 + 1.8 + 0.3 = 9.7\,\%$; gap $9.7 - 5.98 = 3.72$ points.

\textbf{11.} Gaps $3.28$, $2.12$, $2.51$, $2.23$: average about 2.5 points.

\textbf{12.} $16.25 - 2 = 14.25\,\%$.

\textbf{13.} The rule asks for $16.25 - 10.2 = 6.05\,\%$; the actual real rate was $-4.78\,\%$.

\textbf{14.} $2 + 10.2 + 4.1 - 1 = 15.3\,\%$.

\textbf{15.} Inflation and the \omterm{def:g11:money-and-banking:policy-rate}{policy rate} are measured; the neutral real rate and \omterm{def:g12:aggregate-demand-and-supply:potential-output}{potential output} (so the gap) are estimated.

\textbf{16.} No: the record shows what happened under the policy chosen, not what would have happened under the rule. Inflation might have fallen faster, with a deeper
slump, or for other reasons, such as costs falling back.

\textbf{17.} Its rate affects inflation only after a year or more; if it expects inflation to fall by then, it sets the rate for that future inflation.

\textbf{18.} If wage and price setters see it tolerate high inflation, they may expect more, and bringing inflation down later then costs more output.

\textbf{19.} They were chosen for one country and period; Brenland's economy may respond more or less to the rate, and its shocks may differ.

\textbf{20.} In 2022 the rule's rate exceeded the actual \omterm{def:g11:money-and-banking:policy-rate}{policy rate} by about 10.8 points (16.25 against 5.42\,\%); the rule was reacting to inflation of 10.2\,\%.
\end{solution}
